% Apêndice do exemplo: guia rápido da personalização abntex2-ufv. Apague-o no seu trabalho.

\chapter{Guia rápido do \texttt{abntex2-ufv}}\label{ap:classe}

A classe \texttt{abntex2-ufv} personaliza a classe abn\TeX2 \cite{abntex2classe} para o
manual de normalização da UFV \cite{manualUFV}. O uso mais simples é
\texttt{\textbackslash documentclass\{abntex2-ufv\}}: os padrões já seguem o manual (12pt,
A4, impressão só no anverso, português do Brasil) e a classe cuida do layout (margens, fonte,
títulos, capa, folhas, listas, siglas com \texttt{acro}, legendas e os ajustes do
\texttt{biblatex-abnt}). Ela carrega só o que ela mesma precisa (\texttt{fontenc} e
\texttt{indentfirst}); os pacotes de uso do autor ficam no preâmbulo do documento, onde você os
vê e pode ajustá-los: \texttt{graphicx}, \texttt{csquotes}, \texttt{microtype}, o próprio
\texttt{biblatex} com o estilo \texttt{abnt}, e os de tabelas, como \texttt{booktabs},
\texttt{tabularx} e \texttt{multirow}. O template traz o conjunto mínimo. Para mudar algum
padrão da classe, use as opções do \autoref{qua:opcoes} dentro dos colchetes, como em
\texttt{\textbackslash documentclass[twoside, fonte=arial]\{abntex2-ufv\}}.
Os metadados do \ac{pdf} (título e autor) e as cores dos links (pretos, com citações em azul)
também são definidos pela classe; use \texttt{\textbackslash hypersetup\{...\}} para outra
configuração.

\begin{quadro}[htbp]
\begin{ilustracao}{\caption[Opções da classe abntex2-ufv]{Opções da classe \texttt{abntex2-ufv}}\label{qua:opcoes}}{\elaboracaopropria}
  \renewcommand{\arraystretch}{1.2}
  \small
  \begin{tabularx}{\textwidth}{|>{\ttfamily\raggedright\arraybackslash}p{3.4cm}|X|}
    \hline
    \textnormal{\textbf{Opção}} & \textbf{Efeito} \\
    \hline
    fonte=times & Times 12 (TeX Gyre Termes); é o padrão \\
    \hline
    fonte=arial & Arial 12 (TeX Gyre Heros, equivalente métrico do Helvetica/Arial) \\
    \hline
    fonte=cm & Mantém a fonte do documento (Computer Modern ou a que você carregar) \\
    \hline
    capacampus & Imprime \emph{campus} e curso na capa (o manual não os exige) \\
    \hline
    twoside & Impressão frente e verso, com margens espelhadas (padrão: \texttt{oneside}) \\
    \hline
    openany & Capítulos começam em qualquer página (padrão: \texttt{openright}) \\
    \hline
    11pt, 10pt & Tamanho da letra (padrão e exigido pelo manual: 12pt) \\
    \hline
  \end{tabularx}
\end{ilustracao}
\end{quadro}

\section{Migração do pacote para a classe}\label{sec:migracao}

As versões 1.5 e seguintes, até a 2.0, são versões de \emph{transição}. Até a versão 1.3.1
o uso era a classe \texttt{abntex2} com o pacote \texttt{abntex2-UFV}; a partir da 1.5 a forma
recomendada é a classe \texttt{abntex2-ufv}. O pacote (\texttt{abntex2-ufv.sty}) continua
funcionando, mas emite o aviso ``Uso obsoleto'' e \textbf{será descontinuado na versão 2.0}.
Trabalhos novos devem começar com a classe, e os trabalhos em andamento devem ser migrados
agora, pelos passos abaixo.

\textbf{Atenção (versão 1.7).} Nas versões 1.5 e 1.6 a classe carregava, além do seu próprio
layout, os pacotes \texttt{graphicx}, \texttt{csquotes}, \texttt{microtype},
\texttt{booktabs}, \texttt{tabularx}, \texttt{multirow} e \texttt{biblatex}. A partir da 1.7
a classe \emph{não} os carrega mais: eles ficam no preâmbulo do documento (como no template),
onde o editor os reconhece e você controla as opções. Quem migrou na 1.5 ou 1.6 e atualiza
para a 1.7 deve acrescentar esses pacotes (veja o exemplo ao final desta seção). A opção
\texttt{nobiblatex} deixou de existir; se for usada, é ignorada com um aviso.

\begin{enumerate}
  \item \textbf{Atualize os arquivos.} Baixe os arquivos da versão 1.7 ou posterior e substitua os antigos.
        O nome do pacote passou a ser em minúsculas: apague o \texttt{abntex2-UFV.sty} antigo
        e use o \texttt{abntex2-ufv.cls} e o \texttt{abntex2-ufv.sty} novos, na mesma pasta do
        seu arquivo principal (no Overleaf, envie os dois arquivos ao projeto).
  \item \textbf{Troque o \texttt{\textbackslash documentclass}.} A classe já aplica os padrões
        do manual, então as opções \texttt{12pt}, \texttt{a4paper}, \texttt{oneside},
        \texttt{openright}, \texttt{english} e \texttt{brazilian} podem ser removidas. Mantenha
        só o que fugir do padrão (\texttt{twoside}, \texttt{openany}, \texttt{11pt},
        \texttt{fonte=arial}, \texttt{capacampus}), que agora vão entre os colchetes da classe.
        As opções que antes iam em \texttt{\textbackslash usepackage[...]\{abntex2-UFV\}}
        também passam para o \texttt{\textbackslash documentclass}.
  \item \textbf{Apague o \texttt{\textbackslash usepackage\{abntex2-UFV\}}}, o
        \texttt{fontenc} e o \texttt{indentfirst}: a classe já os carrega. Mantenha os
        demais pacotes do seu trabalho.
  \item \sloppy\textbf{Mantenha o \texttt{biblatex} e o \texttt{xurl}} no preâmbulo, com as
        opções do template. O \texttt{\textbackslash urlstyle\{same\}} e o
        \texttt{\textbackslash setlength\{\textbackslash emergencystretch\}} podem ser apagados:
        a classe os define. O \texttt{\textbackslash setlength\{\textbackslash bibparsep\}\{10pt\}}
        do modelo antigo também deve ser apagado: ele deixava o espaço entre as referências
        maior que uma linha em branco (o \texttt{biblatex-abnt} já separa as referências por
        uma linha, e o padrão de \texttt{\textbackslash bibparsep} é 0pt). Para outro estilo de
        citação, altere as opções do \texttt{biblatex} no próprio preâmbulo.
  \item \textbf{Apague o filtro do aviso do \texttt{babel}} (as duas linhas com
        \texttt{silence} e \texttt{\textbackslash WarningFilter}, antes do
        \texttt{\textbackslash documentclass}): a classe já faz isso.
  \item \textbf{Mova as imagens para a pasta \texttt{imagens/}} e declare no preâmbulo
        \texttt{\textbackslash graphicspath\allowbreak\{\{imagens/\}\}}; assim,
        \texttt{\textbackslash includegraphics\{figura\}} basta, sem o nome da pasta. Os
        caminhos antigos, como \texttt{imagens/figura} ou arquivos na pasta principal,
        continuam funcionando.
  \item \textbf{Apague o \texttt{\textbackslash hypersetup} e o \texttt{\textbackslash frenchspacing}}
        se forem os do modelo antigo: a classe já define links pretos, citações em azul,
        o título e o autor do \ac{pdf} e o espaçamento após ponto. Mantenha o
        \texttt{\textbackslash hypersetup} só se quiser algo diferente.
  \item \textbf{Compile e confira.} Rode o \texttt{latexmk} (ou o seu editor) e verifique
        que o aviso ``Uso obsoleto'' não aparece mais e que as listas, as siglas e as
        referências saíram como antes.
\end{enumerate}

O exemplo abaixo resume as mudanças no preâmbulo. Antes (versão 1.3.1):

\begin{lstlisting}[style=exemplo, language={[LaTeX]TeX}]
\documentclass[12pt, openright, oneside, a4paper, english, brazilian]{abntex2}
\usepackage{abntex2-UFV}
\usepackage[T1]{fontenc}
\usepackage{indentfirst}
\usepackage{graphicx}
\usepackage{csquotes}
\usepackage[final]{microtype}
\usepackage{booktabs}
\usepackage{tabularx}
\usepackage{multirow}
\usepackage[backend=biber, style=abnt, uniquename=init, giveninits,
            repeatfields, justify, noslsn]{biblatex}
\usepackage{xurl}
\urlstyle{same}
\addbibresource{referencias.bib}
\setlength\bibparsep{10pt}
\setlength{\emergencystretch}{3em}
\hypersetup{colorlinks=true, linkcolor=black, citecolor=blue, urlcolor=blue}
\end{lstlisting}
Depois (versão 1.7 ou posterior):

\begin{lstlisting}[style=exemplo, language={[LaTeX]TeX}]
\documentclass{abntex2-ufv}
\usepackage{graphicx}
\graphicspath{{imagens/}}
\usepackage{csquotes}
\usepackage[final]{microtype}
\usepackage[backend=biber, style=abnt, uniquename=init, giveninits,
            repeatfields, justify, noslsn]{biblatex}
\usepackage{xurl}
\addbibresource{referencias.bib}
\end{lstlisting}

\section{Dados do trabalho}

Os dados são declarados no preâmbulo e impressos na capa, na folha de rosto e na folha
de aprovação (\autoref{qua:dados}).

\begin{quadro}[htbp]
\begin{ilustracao}{\caption{Comandos de dados do trabalho}\label{qua:dados}}{\elaboracaopropria}
  \renewcommand{\arraystretch}{1.2}
  \small
  \begin{tabularx}{\textwidth}{|>{\ttfamily\raggedright\arraybackslash}p{6.4cm}|X|}
    \hline
    \textnormal{\textbf{Comando}} & \textbf{Função} \\
    \hline
    \textbackslash titulo\{...\}, \textbackslash autor\{...\} & Título e autor \\
    \hline
    \textbackslash local\{...\}, \textbackslash data\{...\} & Cidade e ano de depósito (\texttt{\textbackslash the\textbackslash year} usa o ano atual) \\
    \hline
    \textbackslash instituicao\{...\} & Nome da instituição \\
    \hline
    \textbackslash campus\{...\}, \textbackslash curso\{...\} & \emph{Campus} e curso (capa só com a opção \texttt{capacampus}; o curso também entra na nota descritiva) \\
    \hline
    \textbackslash orientador[sigla]\{...\} & Orientador; \texttt{\textbackslash orientadora} no feminino. A sigla é a da instituição (padrão UFV-CRP) \\
    \hline
    \textbackslash coorientador\{...\} & Coorientador; \texttt{\textbackslash coorientadora} no feminino \\
    \hline
    \textbackslash membrobanca[sigla]\{Titulação Nome\} & Membro da banca; repita para cada membro (ordem de impressão = ordem de declaração) \\
    \hline
    \textbackslash databanca\{...\} & Data da defesa; se omitida, a folha de aprovação traz espaço para preencher à mão \\
    \hline
    \textbackslash siglainstituicaopadrao\{...\} & Troca a sigla padrão (UFV-CRP) dos membros da banca \\
    \hline
    \textbackslash preambuloUFV\{tipo\}\{curso\} & Nota descritiva padrão; \texttt{tipo} = \texttt{tcc}, \texttt{bacharelado} ou \texttt{licenciatura} \\
    \hline
    \textbackslash preambulo\{...\} & Nota descritiva livre, para outros tipos de trabalho \\
    \hline
  \end{tabularx}
\end{ilustracao}
\end{quadro}

\section{Elementos pré-textuais}

Imprimem-se na ordem: \texttt{\textbackslash imprimircapa}, \texttt{\textbackslash imprimirfolhaderosto},
(ficha catalográfica, apenas em teses e dissertações),
\texttt{\textbackslash imprimirfolhadeaprovacao}, ambientes \texttt{dedicatoria},
\texttt{agradecimentos} e \texttt{epigrafe}, resumos (\texttt{resumo}, com
\texttt{\textbackslash palavraschave\{a, b, c\}} e \texttt{\textbackslash keywords\{...\}}),
listas de ilustrações, de tabelas, de siglas e de símbolos e, por último, o sumário.
Dedicatória e epígrafe não têm título e ficam na parte inferior da folha, alinhadas à
direita. Resumo: parágrafo único, de 150 a 500 palavras; as palavras-chave vêm separadas
por ponto e vírgula e terminam em ponto (a macro faz isso a partir da lista com vírgulas).

\section{Ilustrações, quadros e códigos}

\begin{itemize}
  \item \texttt{figure} (``Figura''), \texttt{table} (``Tabela''), \texttt{quadro}
        (``Quadro''), \texttt{grafico} (``Gráfico'') e \texttt{codigo} (``Código''):
        cada tipo tem numeração e lista próprias
        (\texttt{\textbackslash listoffigures*}, \texttt{\textbackslash listoftables*},
        \texttt{\textbackslash listofquadros*}, \texttt{\textbackslash listofgraficos*},
        \texttt{\textbackslash listofcodigos*}).
  \item \texttt{\textbackslash elaboracaopropria} imprime ``Fonte: Elaborado pelo próprio autor'';
        \texttt{\textbackslash fonte\{texto\}} imprime outra fonte; \texttt{\textbackslash legenda\{texto\}}
        imprime a legenda (explicação de símbolos, cores ou siglas) e \texttt{\textbackslash nota\{texto\}}
        imprime uma nota. O rodapé fica alinhado à esquerda, em espaço simples e tamanho 10.
  \item \texttt{ilustracao} e \texttt{tabela}: ambientes de sintaxe
        \texttt{\textbackslash begin\{ilustracao\}\allowbreak\{\textbackslash caption\{...\}\}\allowbreak\{rodapé\}}.
        Medem o conteúdo e dão ao título e ao rodapé a largura dele. \texttt{ilustracao}
        centraliza o título (figuras, quadros, gráficos e códigos) e \texttt{tabela} o
        alinha à esquerda, com o conteúdo em tamanho 10 (padrão do \ac{ibge}). Veja o
        Apêndice~\ref{ap:tabelas}.
  \item \texttt{citacao}: citação direta longa, com recuo de 4 cm, fonte menor e espaço simples.
  \item \texttt{alineas} e \texttt{incisos}: listas da NBR 6024.
\end{itemize}

\section{Siglas}

As siglas usam o pacote \texttt{acro}. Declare-as no preâmbulo e use \texttt{\textbackslash ac}
no texto; a lista é impressa por \texttt{\textbackslash listadesiglas}, em ordem alfabética
e só com as siglas efetivamente usadas (\autoref{qua:siglas}). O manual recomenda uma
lista própria para cada tipo: as abreviaturas (p., ed., art.) usam a etiqueta
\texttt{tag = abreviatura}, são citadas com \texttt{\textbackslash acs} e saem na
\texttt{\textbackslash listadeabreviaturas}. Não use \texttt{\textbackslash ac} em resumos,
títulos e legendas: a expansão da primeira ocorrência deve ficar no texto corrido, onde a
sigla aparece pela primeira vez.

\begin{quadro}[htbp]
\begin{ilustracao}{\caption[Comandos de siglas (pacote acro)]{Comandos de siglas (pacote \texttt{acro})}\label{qua:siglas}}{\elaboracaopropria}
  \renewcommand{\arraystretch}{1.2}
  \small
  \begin{tabularx}{\textwidth}{|>{\ttfamily\raggedright\arraybackslash}p{6.8cm}|X|}
    \hline
    \textnormal{\textbf{Comando}} & \textbf{Resultado} \\
    \hline
    \textbackslash DeclareAcronym\{api\}\{short = API, long = Interface de Programação de Aplicações\} & Declara a sigla (no preâmbulo) \\
    \hline
    \textbackslash ac\{api\} & 1ª vez: Interface de Programação de Aplicações (API); depois: API \\
    \hline
    \textbackslash acp\{api\} & Plural: APIs \\
    \hline
    \textbackslash acs\{api\} / \textbackslash acl\{api\} & Só a sigla / só o nome por extenso \\
    \hline
    \textbackslash listadesiglas & Imprime a Lista de siglas \\
    \hline
    \textbackslash DeclareAcronym\{art\}\{short = art., long = artigo, tag = abreviatura\} & Declara uma abreviatura (lista própria) \\
    \hline
    \textbackslash listadeabreviaturas & Imprime a Lista de abreviaturas \\
    \hline
  \end{tabularx}
\end{ilustracao}
\end{quadro}

\section{Glossário}

O glossário é um elemento pós-textual opcional, com termos técnicos e suas definições. O
ambiente \texttt{glossario} imprime o título ``GLOSSÁRIO'' (que entra no sumário) e cada
\texttt{\textbackslash termo\{termo\}\{definição\}} vira um parágrafo em espaço simples,
com o termo em negrito. Veja o exemplo antes dos apêndices deste modelo.

\section{Compilação}

\begin{itemize}
  \item \textbf{Linha de comando (MiKTeX ou TeX Live)}: na pasta do seu trabalho,
        \texttt{latexmk} executa \texttt{pdflatex} e \texttt{biber} quantas vezes forem
        necessárias. O arquivo \texttt{latexmkrc} grava o \ac{pdf} na pasta do trabalho e
        todos os arquivos auxiliares em \texttt{build/}, e procura o
        \texttt{abntex2-ufv.cls} (e o \texttt{.sty}) em \texttt{../../tex}; se você instalou o
        pacote, ou copiou o template para outro lugar com esses dois arquivos na mesma pasta, apague essa
        linha do \texttt{latexmkrc}.
  \item \textbf{Onde começar}: a pasta \texttt{doc/template} traz um trabalho vazio, pronto
        para ser preenchido; este documento (\texttt{doc/abntex2-ufv.pdf}, com o trabalho de
        exemplo em \texttt{doc/example}) é a documentação.
  \item \textbf{Editores}: TeXstudio e VS Code (LaTeX Workshop) trazem receitas com
        \texttt{latexmk}; configure o motor como \texttt{pdfLaTeX} e a bibliografia como \texttt{biber}.
  \item \textbf{Overleaf}: compilador \texttt{pdfLaTeX}; o \texttt{biber} é executado
        automaticamente.
  \item \textbf{Arquivos auxiliares}: se o PDF ficar com ``(?)'' ou com o sumário
        desatualizado, apague os arquivos \texttt{.aux}, \texttt{.toc}, \texttt{.bbl}
        etc. e compile de novo.
\end{itemize}
